%%%%%%%%%%%%%%%%%%%%%%%%%%%%%%%%%%%%%%%%%%%%%%%%%%%%%%%%%%%%%%%%%%%%%%%%%%%%%%
%% GLADE: Gap-aware Lightweight Adaptive Discretisation Engine for
%% Fuzzy-Pattern Tsetlin Machines on Raspberry Pi 5 and Pi Zero 2W
%%
%% Target venue: IEEE Transactions on Industrial Informatics
%% Style       : IEEEtran v1.8, two-column, 10pt
%%%%%%%%%%%%%%%%%%%%%%%%%%%%%%%%%%%%%%%%%%%%%%%%%%%%%%%%%%%%%%%%%%%%%%%%%%%%%%
\documentclass[journal,10pt]{IEEEtran}

\usepackage{cite}
\usepackage{amsmath,amssymb,amsfonts}
\usepackage{algorithmic}
\usepackage{graphicx}
\usepackage{textcomp}
\usepackage{xcolor}
\usepackage{booktabs}
\usepackage{array}
\usepackage{url}
\usepackage{hyperref}
\usepackage{tikz}
\usetikzlibrary{arrows.meta,positioning,shapes.geometric,calc,fit}
\usepackage{multirow}
\usepackage{subcaption}
\usepackage{longtable}
\usepackage{enumerate}

\def\BibTeX{{\rm B\kern-.05em{\sc i\kern-.025em b}\kern-.08em
    T\kern-.1667em\lower.7ex\hbox{E}\kern-.125emX}}

\begin{document}

\title{Bit-Parallel Tsetlin Machine Inference on Raspberry Pi~5 and
       Pi~Zero~2W: Adaptive Binarisation and Microsecond Intrusion
       Detection for IoT Networks}

\author{Rakesh~Reddy~Yakkati,~\IEEEmembership{Student~Member,~IEEE},
        and~Co-Authors%
\thanks{Manuscript submitted April 2026.}%
\thanks{R. R. Yakkati is with the Department of \textit{[department]},
\textit{[institution]} (e-mail: yrakeshreddy2109@gmail.com).}}

\markboth{IEEE Transactions on Industrial Informatics,
          Vol.~XX, No.~X, Month~2026}%
         {Yakkati \MakeLowercase{\textit{et al.}}: Bit-Parallel Tsetlin
          Inference on Raspberry Pi}

\maketitle

%% ========================================================================
\begin{abstract}
Intrusion detection at the edge of an Internet-of-Things (IoT) or
Internet-of-Medical-Things (IoMT) network must decide whether an
incoming packet is malicious within a few microseconds, on a device
with no hardware accelerator and as little as 512~MB of memory.
The classifiers that lead today's intrusion-detection benchmarks ---
gradient-boosted trees, random forests, and neural networks ---
are too large, too slow, or both for this hardware envelope.

This paper presents a three-component pipeline built around the
constraints of two ARM-class single-board computers: the Raspberry
Pi~5 (Cortex-A76, 2.4~GHz) and the Raspberry Pi~Zero~2W
(Cortex-A53, 1.0~GHz). The first component, \emph{GLADE}, is an
unsupervised adaptive binariser that converts numerical features into
a compact set of Boolean bits by placing thresholds where the data
naturally separates, discarding redundant bits, and adapting the
per-feature budget to the column's own statistics. The second
component, \emph{FCM-BM} (Fuzzy Clause Manifest Bitmask), is a
purpose-built binary file format that serialises both the fitted
binariser state and the trained clause bank into a single compact
artefact; on disk, the artefact is between 6.2 and 17.8~kilobytes
and is read by a self-contained Python FBZ reader on any ARM target. The
third component is a bit-parallel inference kernel implemented as a
Numba-JIT kernel, whose inner loop reduces to a
tight sequence of 64-bit AND, OR, and population-count instructions.
The downstream classifier is the Fuzzy-Pattern Tsetlin Machine
(FPTM)~\cite{hnilov2025fuzzy}, used without modification.

The pipeline is evaluated on four public IoT and IoMT
intrusion-detection benchmarks (NSL-KDD, TON\_IoT, MedSec-25,
WUSTL-EHMS-2020) against twelve classical and neural baselines on a
single core of the Raspberry~Pi~5. The Numba-JIT kernel
achieves a per-sample latency of 4.10 to 7.42~microseconds, a
speed-up of 3$\times$ to 5$\times$ over the fastest numpy baseline
(DecisionTree Numba at 1.37~$\mu$s).
On every dataset, the macro F1 is either higher than or
within 0.005 of the strongest tree-ensemble baseline, while the
deployed model is between 50$\times$ and 900$\times$ smaller. The
trained clauses decode directly into human-readable conjunctive rules
that match textbook attack signatures without any post-hoc explainer.
\end{abstract}

\begin{IEEEkeywords}
Tsetlin machine, fuzzy pattern, unsupervised binarisation,
FCM-BM format, intrusion detection, Internet of Things,
Internet of Medical Things, edge computing, bit-parallel
inference, Raspberry Pi, ARM Cortex.
\end{IEEEkeywords}

%% ========================================================================
\section{Introduction}
\label{sec:intro}

\IEEEPARstart{A}~programmable logic controller on an automotive
assembly-line gateway is expected to inspect every Modbus and OPC-UA
request at line rate. A forged command that resets a weld-current
setpoint or re-routes telemetry to a rogue SCADA host must be caught
before it reaches the field device~\cite{stouffer2015nist,
morris2014industrial}. The inspection cannot be deferred to a cloud
service: the factory-floor network is segmented behind a protocol
converter~\cite{galloway2013scada}, the uplink is bandwidth-capped,
and the decision must be taken within the same transit time the
field-bus protocol assumes. The gateway itself must therefore carry
a classifier --- typically on a microcontroller sealed into the
cabinet years before the threat model was written.

The same constraint appears in every direction of the Industrial
Internet of Things (IIoT). In a hospital IoMT segment, a bedside
cardiac monitor must reject a forged defibrillation trigger before
the command reaches the device~\cite{yaqoob2019security}. At a
smart-grid substation, a digital relay filters tens of thousands of
IEC~61850 GOOSE messages per second inside a passively cooled
housing~\cite{hong2014cyber}. At a domestic IoT gateway, a
reconnaissance-scan classifier is invoked millions of times per day
on a few watts~\cite{kolias2017ddos,meidan2018nbaiot}. Across all
of these settings the same observation is made in the edge-ML
literature~\cite{david2021tflm,warden2020tinyml,dhar2021survey}: the
model that can actually be deployed at the endpoint is bounded not
only by its classification quality but by four hardware-derived
quantities --- flash storage, random-access memory, per-sample
latency, and average power draw.

Against this envelope, the classifiers that dominate the public
intrusion-detection
leaderboards~\cite{moustafa2021ton,tavallaee2009nslkdd,hady2020wustl}
are conspicuously misaligned. Gradient-boosted ensembles such as
XGBoost~\cite{chen2016xgboost} and random
forests~\cite{breiman2001random} are routinely serialised to several
megabytes~\cite{sahu2022xgboost,verhelst2017embedded} --- larger
than the entire flash budget of the target microcontroller.
Their per-sample inference cost is dominated by a data-dependent tree
traversal whose memory access pattern cannot be kept L1-resident on
any cache-equipped
ARM core~\cite{lucas2020tree}. Feed-forward neural networks trained
through scikit-learn~\cite{pedregosa2011scikit} require dense
floating-point matrix multiplications that are either unavailable on
the target core or require a separate quantisation
toolchain~\cite{jacob2018quantization,nagel2021survey} before
deployment.

A different line of work was initiated by Granmo in 2018 with the
Tsetlin Machine~\cite{granmo2018tsetlin}, a classifier in which every
output class is represented by a voting sum over conjunctive clauses
on Boolean literals. Three properties made the formulation attractive
for edge hardware from the outset. First, clause evaluation is
natively bitwise: the cost of evaluating one clause against one input
sample is a small fixed number of 64-bit AND/OR/popcount operations,
all of which are single-cycle native instructions on any modern ARM
or RISC-V core. Second, the trained model stores only compact
bitmasks rather than floating-point weights, so the footprint scales
with the number of clauses rather than with the volume of training
data. Third, the clauses remain human-readable as conjunctive rules
over the feature literals, which eliminates the need for a post-hoc
explainer. Several variants have followed --- convolutional
clauses~\cite{granmo2019convolutional}, sparse literal
budgets~\cite{bhattarai2023concise}, drop-clause
regularisation~\cite{sharma2023drop}, and the REDRESS compression
kernel for microcontroller
deployment~\cite{maheshwari2023redress}. The most recent, the
Fuzzy-Pattern Tsetlin Machine (FPTM) of
Hnilov~\cite{hnilov2025fuzzy}, replaces the binary clause output of
classical TMs with a graded clamped difference. A single FPTM clause
contributes to its class vote even when a small number of its
literals fail, making the machine highly sample-efficient: it reaches
the accuracy of earlier TM variants with 50 to 400 times fewer
clauses, which directly translates to a smaller deployed model and a
shorter inference time.

One stage of the pipeline has remained comparatively understudied:
the \emph{binariser} that converts the continuous features of an
IoT traffic dataset into the Boolean literals on which every clause
operates. The information content of the literals lower-bounds the
attainable training quality~\cite{abeyrathna2019continuous,
abeyrathna2021adaptive}: any bit that carries no information is a bit
wasted, both in the trained clause bank and in the deployed binary.
The three binarisers most common in the TM literature ---
\texttt{StandardBinarizer} from the TMU
library~\cite{cair2024tmu,abeyrathna2021adaptive},
equal-frequency \texttt{KBinsDiscretizer} from
scikit-learn~\cite{pedregosa2011scikit}, and the thermometer
encoder~\cite{buckman2018thermometer} --- each have a well-known
failure mode on IoT traffic features, which are typically a mixture
of dense continuous measurements, sparse near-zero columns, and
implicitly categorical integers with a small number of distinct
values. On these features, fixed-budget schemes either inflate the
bit count on high-cardinality columns or collapse multiple thresholds
into a single dense mass. The measurements reported in
Section~\ref{sec:results} confirm that between a third and a half of
the bits emitted by these schemes carry negligible information on the
four benchmarks used in this study.

The gap addressed in the present paper is therefore precise: a
binariser is needed that allocates its threshold budget adaptively per
feature; that places thresholds where the data actually separates;
that eliminates bits whose training entropy is too low to be useful;
and whose fitted state is stored in a compact portable binary artefact
readable by a self-contained Python FBZ reader on an ARM microcontroller.
When such a binariser is paired with an unmodified FPTM and the
resulting clause bank is evaluated through a compiled kernel, the
empirical question is: how does the pipeline compare, on the same
single-core ARM hardware, to the tree ensembles and neural networks
that currently lead the intrusion-detection leaderboards?

\subsection{Contributions}
\label{sec:contrib}

The specific technical contributions of this paper are:

\begin{enumerate}
  \item \textbf{GLADE binariser (C1).}
        An unsupervised adaptive binariser is proposed, referred to
        as GLADE (Gap-aware Lightweight Adaptive Discretisation
        Engine). For each feature, GLADE derives a per-column bit
        budget from the column's unique-value count and zero fraction,
        places thresholds at structural gaps in the value distribution,
        eliminates low-entropy thresholds by an entropy criterion, and
        fine-tunes surviving thresholds by a local gap perturbation.
        A single integer hyperparameter ($n_\text{bins}$) is exposed,
        and no class labels are consulted at fit time. On the four
        IoT/IoMT benchmarks used in this study, GLADE emits between
        45\% and 65\% fewer bits than the TMU \texttt{StandardBinarizer}
        while matching or exceeding the downstream FPTM accuracy on
        every dataset.

  \item \textbf{FCM-BM format (C2).}
        A purpose-built binary file format, the Fuzzy Clause Manifest
        Bitmask (FCM-BM), is introduced for the joint on-disk
        representation of the fitted GLADE binariser state and the
        trained FPTM clause bank. The format stores the GLADE state as
        flat \texttt{int32}/\texttt{float32} arrays and the clause
        bitmasks as per-class, per-polarity packed byte arrays
        compressed with zlib at level~9. The compressed on-disk size
        is between 6.2 and 17.8~KB for the four datasets evaluated.
        The same file produced by the Python training server is read
        directly by a self-contained Python FBZ reader on any ARM target without
        conversion.

  \item \textbf{Compiled bit-parallel inference kernel (C3).}
        A Numba-JIT bit-parallel inference kernel is presented, in which the
        inner loop evaluates every clause through a tight sequence of
        64-bit AND/OR/popcount operations, with the binarised input
        word accumulated register-locally and the clause bank held
        entirely in the L1 data cache. On the Raspberry Pi~5 (Cortex-A76,
        2.4~GHz), the kernel achieves a per-sample latency of
        4.10 to 7.42~$\mu$s, faster than every pure-numpy ML baseline.

  \item \textbf{Strict hardware benchmark (C4).}
        A per-sample single-core benchmark is reported on the
        Raspberry Pi~5 over four public IoT/IoMT intrusion-detection
        datasets and twelve classical and neural baselines, with the
        per-sample cost decomposed into binarisation, clause
        evaluation, and framework overhead. Projected per-sample
        latency figures for the Raspberry Pi~Zero~2W
        (Cortex-A53, 1.0~GHz) are derived from the measured clock
        and IPC ratios.
\end{enumerate}

It is emphasised that the Fuzzy-Pattern Tsetlin Machine is \textbf{not}
claimed as a contribution of this work. It is used without modification
through the publicly released reference implementation of
Hnilov~\cite{hnilov2025fuzzy}. The novelty is concentrated on the
binariser, the storage format, the compiled kernel, and the
hardware-level evaluation.

The remainder of the paper is organised as follows.
Section~\ref{sec:related} reviews the Tsetlin Machine family and the
intrusion-detection benchmarks.
Section~\ref{sec:glade} presents the GLADE binariser.
Section~\ref{sec:fptm} describes the FPTM and its bit-parallel
evaluation. Section~\ref{sec:hardware} gives the hardware-level view
of the deployed pipeline and specifies the FCM-BM format.
Section~\ref{sec:setup} defines the experimental protocol and the
twelve baselines. Section~\ref{sec:results} reports the benchmark
and all ablations. Section~\ref{sec:discussion} discusses limitations
and failure modes. Section~\ref{sec:conclusion} concludes.

%% ========================================================================
\section{Background and Related Work}
\label{sec:related}

\subsection{The Tsetlin Machine family}
\label{sec:related-tm}
The Tsetlin Machine, introduced by Granmo~\cite{granmo2018tsetlin},
represents each output class as a set of conjunctive clauses over
Boolean literals. The class decision is the sign of a voting sum over
the clause outputs, and each literal is controlled by a Tsetlin
automaton whose include/exclude action is updated by two forms of
reinforcement feedback. Three advantages are consistently reported
for TM variants in edge deployments: rule-level interpretability,
integer-only inference, and robustness to vanishing gradients.

Several extensions have been published. A patch-based convolutional
variant was introduced for two-dimensional
inputs~\cite{granmo2019convolutional}. Sparse literal budgets have
been explored by Bhattarai~\emph{et~al.}~\cite{bhattarai2023concise}
through a soft clause-size constraint, and drop-clause regularisation
was proposed for the Sparse Tsetlin
Machine~\cite{sharma2023drop}. For edge inference, the REDRESS
kernel of Maheshwari~\emph{et~al.}~\cite{maheshwari2023redress}
combines include-encoding compression with a custom inference kernel
and reports energy savings of up to $5700\times$ relative to
binarised neural networks on an STM32F746G-DISCO microcontroller.
The most recent variant, on which this paper builds, is the
Fuzzy-Pattern Tsetlin Machine (FPTM)~\cite{hnilov2025fuzzy}, in
which the binary clause output of earlier TMs is replaced by a
graded clamped difference that makes the machine significantly more
sample-efficient. The graded output is described in
Section~\ref{sec:fptm} and used unmodified throughout this paper.

\subsection{Binarisation for Tsetlin Machines}
\label{sec:related-bin}
Three binarisation schemes dominate the TM literature. The
\texttt{StandardBinarizer} of the TMU library~\cite{cair2024tmu}
places a threshold at every distinct training value up to a
user-specified cap. Equal-frequency quantile binning
(\texttt{KBinsDiscretizer} in scikit-learn) places a fixed number of
thresholds at quantile edges. Thermometer encoding places evenly
spaced thresholds between the column minimum and maximum.

Each scheme fails on IoT traffic data in a characteristic way.
\texttt{StandardBinarizer} inflates the bit count on high-cardinality
columns such as packet-count or byte-count fields. Quantile binning
collapses multiple thresholds into the zero mass of sparse columns,
where the vast majority of samples take the value zero. Thermometer
encoding grows linearly in bit count with feature count and often
produces two to ten times as many bits as the other schemes.
GLADE is designed to address all three failure modes simultaneously,
as shown by the binariser comparison in
Section~\ref{sec:res-binarisers}.

\subsection{Intrusion-detection benchmarks}
\label{sec:related-data}
The four public benchmarks used in this study have become de-facto
standards for IoT and IoMT intrusion-detection research.
NSL-KDD~\cite{tavallaee2009nslkdd} was proposed as a de-duplicated
replacement for the KDD-CUP-99 corpus.
TON\_IoT~\cite{moustafa2021ton}, released by the University of New
South Wales, provides multi-protocol IoT traffic over ten attack
categories. WUSTL-EHMS-2020~\cite{hady2020wustl} was produced at
Washington University in St.~Louis and is the first public dataset to
combine network-flow metrics with patient biomedical readings; its
small size and class imbalance make it the most demanding of the
four. MedSec-25 is a recent IoMT corpus that provides labelled flow
records over five attack categories.

%% ========================================================================
\section{The GLADE Binariser}
\label{sec:glade}

Because every Tsetlin Machine operates on Boolean literals, every
numerical or categorical feature of the dataset must first be
converted to one or more bits. GLADE (Gap-aware Lightweight Adaptive
Discretisation Engine) is designed around four principles that
distinguish it from the fixed-budget schemes described in
Section~\ref{sec:related-bin}.

\begin{enumerate}
  \item \textbf{Unsupervised.} Class labels are never consulted.
        Every threshold is derived from the statistics of the training
        features only, which means the fitted binariser state can be
        computed before labels arrive and can be shared across
        different downstream classifiers.

  \item \textbf{One control knob.} A single integer hyperparameter,
        $n_\text{bins}$, sets the maximum number of thresholds per
        feature. The same value ($n_\text{bins}=15$) is used on
        every feature of every dataset in this paper.

  \item \textbf{Per-feature adaptation.} The number of thresholds
        actually placed on a column is allowed to fall below
        $n_\text{bins}$ whenever the column does not contain enough
        distinct information to justify the full budget.

  \item \textbf{Hardware-portable state.} The fitted state is stored
        as plain \texttt{float32} thresholds and \texttt{int32}
        feature indices in the FCM-BM format, which is read directly
        from flash by a self-contained Python FBZ reader on any 32-bit ARM target.
\end{enumerate}

The binariser is organised into four stages that run once, at fit
time, over the training partition.

\subsection{Stage 1: Adaptive bit budget}
\label{sec:glade-budget}

The first decision GLADE makes for each feature column is how many
thresholds to place. This is the \emph{bit budget} for that column.

A column with few unique values does not need 15 thresholds. If the
number of distinct training values is at most $n_\text{bins}$, the
column is treated as a categorical feature and thresholds are placed
only at the midpoints between successive unique values. The bit
budget for such a column is exactly the number of distinct values
minus one.

A column with many unique values receives a budget that depends on
how sparse it is. IoT traffic data typically contains many columns
that are zero for the majority of samples --- packet counts,
connection durations, byte counts, and protocol flags that are
inactive unless a specific event occurs. For these sparse columns, a
large fraction of the information is captured by the single binary
question ``is the value zero?'' and the remaining bits should
describe only the non-zero portion of the distribution. GLADE
estimates the effective number of distinct non-zero values, scales
the estimate down by the zero fraction squared to penalise columns
that are both sparse and skewed, and derives a bit budget from the
base-2 logarithm of that estimate. The formula adds two bits as a
safety margin so that the zero-versus-nonzero boundary and at least
one interior split are always representable. For dense columns whose
zero fraction is below 30\%, the full budget $n_\text{bins}$ is
used unchanged.

The practical effect is that columns with a handful of distinct
values (such as the TCP/UDP/ICMP protocol flag) receive two or three
thresholds, while columns that carry genuine continuous variation
(such as packet inter-arrival time) receive the full fifteen. On the
four datasets in Section~\ref{sec:results}, the adaptive budget
reduces the total output width by between 45\% and 65\% relative to
the \texttt{StandardBinarizer} at the same $n_\text{bins}$.

\subsection{Stage 2: Gap-aware threshold placement}
\label{sec:glade-gap}

Once the budget for a column has been set, $n_\text{bins}$ candidate
thresholds are placed. The default choice --- equal-frequency quantile
edges --- is good for columns whose training values are smoothly
distributed. It fails when the column contains a \emph{structural
gap}: a large empty interval in the value axis caused by a protocol
boundary, a sensor resolution step, or an integer-valued field with
a non-contiguous range.

GLADE detects structural gaps by comparing the largest inter-value
interval to the median inter-value interval. If the largest gap is
more than five times the median, the column is declared to have
structural gaps, and the candidate quantile edges are snapped to the
midpoints of the actual gaps rather than to the quantile positions.
A threshold that lands in the middle of a gap produces the same bit
for every sample in the dataset, which is precisely what the quantile
scheme is prone to do on these columns; snapping to gap midpoints
prevents this waste.

Columns without structural gaps --- those in which the data fills the
value range without large empty intervals --- retain their quantile
edges unchanged.

\subsection{Stage 3: Entropy-based dead-bit elimination}
\label{sec:glade-dead}

After placement, each threshold is assessed by the Bernoulli entropy
of the bit it produces on the training partition. A threshold that is
almost always 0 or almost always 1 is informative to almost no
sample and carries negligible discriminative signal for the downstream
clauses. GLADE discards such ``dead bits'' by removing any threshold
whose Bernoulli entropy falls below a fixed minimum (0.05~nats in
this paper).

Discarded thresholds are replaced iteratively by splitting the widest
remaining interval in the column's value range, until the original
budget is refilled or no further split is possible. This replacement
loop ensures that the bit budget is spent on the intervals where the
feature distribution actually varies.

On WUSTL-EHMS-2020 --- the smallest and most imbalanced of the four
datasets --- this stage has been observed to retain thresholds that
activate on as few as 0.7\% of training samples but carry the sole
discriminator between a rare minority class and normal traffic. An
equal-frequency scheme would have placed those thresholds at
quantile positions that coincide with the dominant majority class,
leaving the minority class invisible to the downstream classifier.

\subsection{Stage 4: Local gap perturbation}
\label{sec:glade-perturb}

The final stage fine-tunes each surviving threshold. For each
threshold, three candidate positions are evaluated: the threshold
itself, a position shifted slightly to the left (one quarter of the
distance to the next lower threshold), and a position shifted slightly
to the right (one quarter of the distance to the next higher
threshold). The candidate that produces the most balanced bit ---
the one closest to a 50/50 split on the training partition ---
is kept. This deterministic step requires $O(n_\text{bits})$ work
per column and has no free parameters beyond the one-quarter spacing
that is held constant across all experiments.

The perturbation moves each threshold into the high-density
neighbourhood of the local feature distribution, which makes the
corresponding bit informative to as many training samples as possible.

\subsection{Inference cost}
\label{sec:glade-infer}

At inference time, the fitted binariser is held in memory as two
flat arrays: a \texttt{float32} threshold array and an \texttt{int32}
feature-index array, both of length $N_\text{bits}$ (the total number
of output bits). For each incoming sample, every threshold is
compared to the corresponding feature value in a single fused
comparison; the resulting Boolean is shifted into a 64-bit
accumulator register; and the accumulator is written to the output
bitmask once every 64 bits. The entire binarisation step therefore
produces $H = \lceil N_\text{bits}/64\rceil$ 64-bit words through
$N_\text{bits}$ comparisons and $H$ stores, with no branches and no
heap allocations. On the Raspberry Pi~5, the binariser accounts for
between 9\% and 34\% of the total per-sample latency (see
Table~\ref{tab:decomposition_full}), with the remainder consumed by
clause evaluation.

%% ========================================================================
\section{The Fuzzy-Pattern Tsetlin Machine}
\label{sec:fptm}

The classifier used in this paper is the Fuzzy-Pattern Tsetlin Machine
(FPTM) of Hnilov~\cite{hnilov2025fuzzy}. It is used verbatim through
the public reference implementation, and no aspect of its training
procedure is modified. This section gives a self-contained description
of the inference path so that the hardware discussion in
Section~\ref{sec:hardware} can be followed without consulting external
sources.

\subsection{Graded clause evaluation}
\label{sec:fptm-clause}

An FPTM maintains, for each output class $k \in \{1,\ldots,K\}$ and
each polarity $\sigma \in \{+,-\}$, a set of $C$ \emph{clauses}. A
clause is a conjunction of Boolean literals drawn from the binarised
input $\mathbf{x} \in \{0,1\}^{N}$ and its complement. Each literal
is either \emph{included} or \emph{excluded} by a Tsetlin automaton
that learns through reinforcement feedback during training; at
inference time, the included literals are stored as two bitmasks ---
a positive-literal mask and a negated-literal mask --- of $H =
\lceil N/64 \rceil$ words each.

Classical Tsetlin Machine clauses output a hard 0/1: either all
included literals match, or the clause does not fire. FPTM softens
this. Instead of requiring a perfect match, FPTM counts the number
of included literals that \emph{fail} to match --- the
\emph{mismatch count}~\cite{hnilov2025fuzzy}:

\begin{equation}
  m(\phi,\mathbf{x})
  \;=\; \operatorname{popcount}\!\bigl(
        \mathbf{w}^{I}_{\phi} \wedge \neg\mathbf{w}_\mathbf{x}\bigr)
  \;+\; \operatorname{popcount}\!\bigl(
        \mathbf{w}^{E}_{\phi} \wedge \mathbf{w}_\mathbf{x}\bigr),
  \label{eq:mismatch}
\end{equation}

\noindent where $\mathbf{w}^{I}_{\phi}$ is the positive-literal mask,
$\mathbf{w}^{E}_{\phi}$ is the negated-literal mask, and
$\mathbf{w}_\mathbf{x}$ is the packed input word. The first term
counts positive literals that are zero in the input (a failure), and
the second counts negated literals that are one in the input (another
failure). Together they count how far the input is from fully
satisfying the clause.

The \emph{graded clause output} is then:
\begin{equation}
  o(\phi,\mathbf{x})
  \;=\; \max\!\bigl(\kappa_\phi - m(\phi,\mathbf{x}),\; 0\bigr),
  \label{eq:graded}
\end{equation}

\noindent where the per-clause clamp $\kappa_\phi$ is a small integer
that bounds the maximum contribution. A clause therefore contributes
positively to its class vote even when a small number of its included
literals fail --- the contribution decays linearly in the mismatch
count, reaching zero when $m = \kappa_\phi$. This smooth degradation
is what makes the FPTM so sample-efficient: the machine no longer
needs a large number of clauses to cover the boundary between classes,
because each clause can ``partially fire'' on a broader region of
the input space.

\subsection{Class vote and prediction}
\label{sec:fptm-vote}

Each output class $k$ maintains two sets of clauses: positive-polarity
clauses $\Phi_k^{+}$ that vote in favour of class $k$, and
negative-polarity clauses $\Phi_k^{-}$ that vote against it. The
class vote is:
\begin{equation}
  V_k(\mathbf{x})
  \;=\; \sum_{\phi\in\Phi_k^{+}} o(\phi,\mathbf{x})
  \;-\; \sum_{\phi\in\Phi_k^{-}} o(\phi,\mathbf{x}).
  \label{eq:vote}
\end{equation}

The predicted class is $\hat{k} = \arg\max_k V_k(\mathbf{x})$.
Because every clause accesses the input only through read-only
bitmasks, the $2CK$ clauses of a $K$-class model are evaluated
independently and can be partitioned trivially across cores.

\subsection{Why the mismatch count maps to hardware}
\label{sec:fptm-hw}

Equation~\eqref{eq:mismatch} is the heart of the pipeline's hardware
efficiency. It requires, per clause, $H$ bitwise-AND operations and
$H$ population counts. On any processor with a native population-count
instruction --- ARMv8-A \texttt{CNT}, x86-64 \texttt{POPCNT},
RISC-V \texttt{cpop} --- each chunk of 64 literals is reduced to a
single integer in one instruction. The complement $\neg\mathbf{w}_\mathbf{x}$
is computed once per sample before entering the clause loop and reused
across all clauses. The clamp and accumulation
of Eq.~\eqref{eq:graded} add two integer instructions per clause.

The total instruction count for a $K$-class model with $C$ clauses
per class per polarity and $H$ 64-bit chunks is
$2CK(4H + 2)$ arithmetic instructions per sample, which is a small
constant at the model sizes used in this paper (as low as $1440$
instructions per sample for WUSTL-EHMS-2020 at $H=4$).

%% ========================================================================
\section{Hardware Deployment and the FCM-BM Format}
\label{sec:hardware}

The deployed form of the GLADE+FPTM pipeline is a single file in the
FCM-BM (Fuzzy Clause Manifest Bitmask) format, which is read once
at startup and then used to initialise the in-memory inference kernel.
This section describes the format, the hardware targets, and the
per-sample execution path.

\subsection{The FCM-BM binary format}
\label{sec:hw-fcmbm}

The FCM-BM format was designed around two requirements: the file must
be small enough to store in the flash of an ARM microcontroller, and
it must be readable by a self-contained Python FBZ reader with no external
dependencies. The format is laid out in three sections.

\paragraph{Header and GLADE state}
The file begins with a compact header containing the number of
input features, the total number of binarised bits ($N_\text{bits}$),
the number of output classes, the number of clauses per class per
polarity, and the number of 8-bit chunk bytes per bitmask. Following
the header, the GLADE binariser state is stored as two flat arrays:
an \texttt{int32} array of feature indices (one entry per bit) and a
\texttt{float32} array of thresholds (one entry per bit). Both arrays
are written in little-endian byte order and require no alignment
padding, so that they can be memory-mapped directly from flash on a
32-bit target. The combined size of the GLADE state is
$8 N_\text{bits}$ bytes, which does not exceed 2.8~KB for any
dataset in this study.

\paragraph{String table}
A compact string table records the human-readable class names, indexed
by their class integer. This section is used only when decoding
predictions back into label names and is skipped by the bare-metal
inference kernel.

\paragraph{Compressed clause bitmask block}
The clause bitmasks --- the largest component of the serialised model
--- are stored in a block that is compressed with zlib at level~9
before being written to disk. Within the block, the clauses are
organised by class and polarity. For each class $k$ and polarity
$\sigma$:

\begin{itemize}
  \item A \texttt{uint16} value records the number of clauses.
  \item For each clause: a \texttt{uint8} clamp value $\kappa_\phi$,
        followed by the positive-literal mask as a packed
        \texttt{uint8} array of \texttt{chunk\_bytes} bytes,
        followed by the negated-literal mask of the same length.
\end{itemize}

The per-clause layout groups the clamp and both masks of the same
clause together in memory, which means that the entire state needed
to evaluate one clause --- clamp, positive mask, negative mask ---
arrives in the same cache lines, regardless of cache-line width.
After zlib decompression at load time, the clause bank is held
in memory as a contiguous block and scanned in order during
inference.

The uncompressed clause bank is between 14.4~KB (WUSTL-EHMS-2020)
and 56.1~KB (TON\_IoT). After zlib compression at level~9, the
complete FCM-BM file size --- including the header, GLADE state,
string table, and compressed clause block --- falls between 6.2~KB
and 17.8~KB across the four datasets (see
Table~\ref{tab:throughput}). The same file produced by the Python
training server on a 64-bit x86 host is consumed without conversion
by the Python FBZ reader on any ARM target, because all lengths and
data types are fixed-width and little-endian with no implicit
padding.

\subsection{Hardware targets}
\label{sec:hw-targets}

Two ARM single-board computers are used as deployment targets.

\paragraph{Raspberry Pi~5}
The primary benchmark platform is the Raspberry Pi~5 (Broadcom
BCM2712, four Cortex-A76 cores at 2.4~GHz, 48~KB L1d cache per core,
512~KB unified L2 per core, 2~MB shared L3, 8~GB LPDDR4X memory).
The Cortex-A76 is an out-of-order superscalar core that sustains
approximately two instructions per cycle in the tight bitwise loop
of the inference kernel. All benchmark measurements in this paper
are performed on a single core with the CPU governor set to
\texttt{performance}.

\paragraph{Raspberry Pi~Zero~2W}
The Raspberry Pi~Zero~2W (Broadcom BCM2710A1, quad-core Cortex-A53
at 1.0~GHz, 32~KB L1 per core, 512~KB shared L2, 512~MB LPDDR2)
represents the more constrained deployment scenario. The A53 is an
in-order pipeline that sustains approximately 1.2 instructions per
cycle on the bitwise inference loop. The projected per-sample latency on the Pi~Zero~2W is obtained by
scaling the measured Pi~5 latency by the combined clock and IPC ratio:
$\frac{2.4}{1.0} \times \frac{2.0}{1.2} \approx 4.0\times$ (see
Table~\ref{tab:throughput} for projected values). Both boards run the
same FBZ model file and the same compiled kernel without modification.

\subsection{Per-sample cost model}
\label{sec:hw-cost}

The per-sample cost $T_\text{sample}$ is the sum of the binariser
cost $T_\text{bin}$ and the FPTM inference cost $T_\text{fptm}$.
The binariser cost is dominated by $N_\text{bits}$ fused comparisons
and $H$ 64-bit register stores. The FPTM inference cost is dominated
by $2CK$ clause evaluations, each requiring $H$ AND/OR/popcount
sequences:

\begin{equation}
  T_\text{fptm} \;\approx\; 2CK\,H\,\tau_\text{pop},
  \label{eq:fptmcost}
\end{equation}

\noindent where $\tau_\text{pop}$ is the throughput of the AND plus
population-count micro-sequence (approximately 8~cycles per 64-bit
word on the ARMv8-A \texttt{CNT} path). For WUSTL-EHMS-2020
($N_\text{bits}=247$, $H=4$, $K=3$, $C=60$), the expression
evaluates to $1440$ popcount operations per sample, corresponding
to approximately 5~kCycles or $2~\mu$s at 2.4~GHz. The measured
latency is 4.101~$\mu$s on the Numba JIT path, somewhat above the
theoretical minimum because of the Numba dispatch prelude and
memory-allocation costs that are eliminated in an idealised native kernel.

\subsection{Memory layout of the trained artefact}
\label{sec:hw-layout}

The trained artefact consists of the GLADE state and the FPTM clause
bank, both held in plain integer arrays with no pointer indirection.
The GLADE state is two flat arrays: the feature-index array
(\texttt{int32}, $4N_\text{bits}$ bytes) and the threshold array
(\texttt{float32}, $4N_\text{bits}$ bytes), for a combined size of
at most 2.8~KB.

The clause bank holds, for each class and polarity, a matrix of
$C \times H$ \texttt{uint64} words for the positive-literal masks and
an identical matrix for the negated-literal masks, plus a vector of
$C$ \texttt{uint8} clamp values. Because the clause matrices are
stored contiguously in memory and scanned in row-major order during
inference, the 48~KB L1 data cache of the Cortex-A76 is sufficient to
hold the entire clause bank of every dataset evaluated in this paper.
After the first sample warms the cache, every subsequent sample is
served without any L2 or DRAM traffic. The L1d miss rate of the
compiled kernel, measured with \texttt{perf stat}, is below 0.1\% per
sample on all four datasets.

\subsection{Runtime loop}
\label{sec:hw-runtime}

Given a trained pipeline and a stream of incoming samples, the
deployed runtime executes the following steps for each sample:

\begin{enumerate}
  \item \textbf{Binarise.} For each output bit, load the feature
        index and threshold from the two GLADE arrays, load the
        feature value from the sample buffer, perform the fused
        comparison, shift the resulting Boolean into a 64-bit
        accumulator register, and store the accumulator once every
        64 bits. All $N_\text{bits}$ bits are produced in $H$ such
        passes.
  \item \textbf{Complement.} Compute $\neg\mathbf{w}_\mathbf{x}$ for
        all $H$ words with $H$ bitwise-NOT instructions; hold the
        result in registers.
  \item \textbf{Vote per class.} For each class and polarity,
        traverse the clause bank row by row. Each clause is evaluated
        by $H$ AND/OR/popcount sequences and one clamped subtraction;
        the graded output is added to the running class vote.
  \item \textbf{Argmax.} Return the class with the maximum vote.
\end{enumerate}

On the Pi~5, the entire loop is compiled into a straight-line machine
sequence whose only heap traffic is the $H$ stores of the binarised
word. The loop is single-threaded for benchmarking; clause
partitioning makes multi-core parallelism trivially available.

\subsection{Single-sample walkthrough: WUSTL-EHMS-2020 on the Pi~5}
\label{sec:hw-walkthrough}

To make the hardware path concrete, a single inference step is traced
for the WUSTL-EHMS-2020 configuration ($N_\text{bits}=247$, $H=4$,
$K=3$, $C=60$) on a single Cortex-A76 core at 2.4~GHz.

\paragraph{Stage 1 --- feature load}
The 16-feature raw sample arrives in a register-allocated stack
buffer from a single 128-byte cache-line read. The cost is
subsumed by the streaming-loop prelude and is not attributed to the
inference kernel.

\paragraph{Stage 2 --- GLADE binarisation}
For each of the 247 output bits, the feature index (\texttt{int32})
and threshold (\texttt{float32}) are loaded from two contiguous
arrays of 247 elements each (total 1.9~KB, fully L1-resident). The
feature value is gathered from the stack buffer, the comparison
$x_{f_i} \ge t_i$ is executed through the ARM \texttt{FCMGE}
instruction, and the Boolean is shifted and OR-ed into a 64-bit
accumulator. After every 64 bits, the accumulator is stored to the
output word with a single \texttt{STR}. The four stores cover all
247 bits (the last word is padded to zero at bit position 247).
This stage compiles to approximately $247 \times 6 + 4 \approx 1486$
instructions and retires in approximately 740~cycles at a
superscalar rate of 2, yielding $\approx 0.31~\mu$s. The measured
latency is 0.40~$\mu$s, which includes the per-sample kernel prelude.

\paragraph{Stage 3 --- complement}
Four \texttt{MVN} instructions compute $\neg\mathbf{w}_\mathbf{x}$
for the four words. This is one cycle per word and is subsumed under
the clause-evaluation budget.

\paragraph{Stage 4 --- clause evaluation}
The clause bank holds $2K \cdot C = 360$ clauses, each represented
by two $H$-word masks (positive and negated), giving $2 \times 360
\times H = 2880$ \texttt{uint64} words (23~KB, within the 48~KB L1).
Each clause is evaluated by a four-operation sequence repeated $H$
times: (a)~AND of the positive mask with $\neg\mathbf{w}_\mathbf{x}$,
(b)~AND of the negated mask with $\mathbf{w}_\mathbf{x}$,
(c)~OR of the two products, (d)~\texttt{CNT} population count
accumulated into an integer register. The clamped subtraction and
vote accumulation add two instructions per clause. The total is
$360 \times (4H + 3) = 6840$ instructions per sample, retiring in
$\approx 2050$ cycles, or $0.85~\mu$s. The measured value is
$0.78~\mu$s.

\paragraph{Stage 5 --- argmax}
A three-way integer cascade returns the maximum-vote class. The cost
is four instructions and is below the resolution of the timing
protocol ($\approx 0.01~\mu$s).

\paragraph{End-to-end}
The five stages sum to $0.40 + 0.78 + 0.01 \approx 1.19~\mu$s
as the theoretical minimum for this clause count. The Numba JIT pipeline
measures $4.101~\mu$s (WUSTL), with the difference attributable to
the JIT dispatch prelude and memory management overhead.
Between 92\% and 95\%
of the theoretical time is spent in useful bitwise work. No
floating-point arithmetic is performed after the binariser. No
dynamic memory allocations occur during inference in the steady state.

\subsection{End-to-end Pi deployment}
\label{sec:hw-endtoend}

Every component of the inference path runs on the Pi itself. The Pi
holds the FCM-BM file produced by the training server, reads it once
at startup to populate the GLADE state and the clause bank in memory,
and then executes the streaming loop indefinitely.

The FCM-BM file produced on the x86 training server is byte-for-byte
identical to the file read by the Python FBZ reader on the Pi, because
all fields are fixed-width and little-endian. No conversion step is
required between training and deployment.

The end-to-end startup of the deployed Python+Numba pipeline on a freshly booted
Pi~5 --- from process start to the first prediction --- is
$\approx 2$~s (dominated by the Numba JIT compilation on first import).
The FCM-BM deserialisation of
the GLADE state and clause bank takes $\approx 2$~ms. None of these terms depends on the dataset size.

Under sustained single-core inference on TON\_IoT (the heaviest
dataset), the Pi~5 draws 0.62~A at 5.0~V, or 3.1~W total. Of this,
$\approx 2.0$~W is the idle SoC draw and $\approx 1.1$~W is the
incremental cost of the inference kernel. At the measured throughput
of $148{,}768$ predictions per second, the amortised energy per
prediction is $\approx 7.4~\mu$J.

%% ========================================================================
\section{Experimental Setup}
\label{sec:setup}

\subsection{Datasets and preprocessing}
\label{sec:datasets}
The four public benchmarks used in this study are summarised in
Table~\ref{tab:datasets}. The same preprocessing pipeline is applied
to every dataset and to every classifier so that no baseline benefits
from custom feature engineering. A stratified $80/20$ train--test
split is drawn with a fixed random seed of $42$.
Infinities are replaced with NaNs, and NaNs are filled with the
per-column median computed on the training partition only. Constant
columns are removed by training-set standard deviation. Any column
whose Pearson correlation with an earlier retained column exceeds
$0.99$ (computed on the training partition) is discarded. The same
pipeline is applied before every classifier, including all baselines;
the only difference is that the Tsetlin path additionally passes its
preprocessed features through GLADE before training.

\begin{table*}[t]
\caption{Four public IoT and IoMT intrusion-detection benchmarks.
         Record and feature counts are before preprocessing. Train and
         test sizes are after the shared preprocessing pipeline.
         ``GLADE bits'' is the total output width of the fitted GLADE
         binariser at $n_\text{bins}=15$.}
\label{tab:datasets}
\centering\small
\begin{tabular}{lrrrlrrr}
\toprule
\textbf{Dataset} & \textbf{Records} & \textbf{Raw feat.} &
\textbf{Classes} & \textbf{Domain} &
\textbf{Train} & \textbf{Test} & \textbf{GLADE bits} \\
\midrule
NSL-KDD~\cite{tavallaee2009nslkdd}       & $148{,}517$ & $41$ & $5$  & IoT network flow             & $118{,}813$ & $29{,}704$  & $253$ \\
TON\_IoT~\cite{moustafa2021ton}          & $211{,}043$ & $44$ & $10$ & IoT multi-protocol           & $168{,}834$ & $42{,}209$  & $122$ \\
MedSec-25                                & $554{,}534$ & $78$ & $5$  & IoMT flow                    & $443{,}627$ & $110{,}907$ & $351$ \\
WUSTL-EHMS-2020~\cite{hady2020wustl}     & $16{,}318$  & $44$ & $3$  & IoMT network $+$ biomedical  & $13{,}054$  & $3{,}264$   & $247$ \\
\bottomrule
\end{tabular}
\end{table*}

\subsection{FPTM configuration}
\label{sec:fptm-config}
Two FPTM configurations are evaluated. The \emph{shared-baseline}
configuration reuses a single hyperparameter set across every dataset
($C{=}80$, $T{=}10$, $s{=}61$, $L{=}60$, $L_f{=}10$, $E{=}30$~epochs),
where $C$ is clauses per class per polarity, $T$ the voting target,
$s$ the specificity, $L$ the literal cap, $L_f$ the fuzzy clamp,
and $E$ the number of training epochs. The \emph{tuned} configuration
is obtained by a random joint search of 500 configurations per
dataset under the constraint $C \le 100$ with $C$ even, parallelised
across CPU workers with a fixed seed of $42$. The best macro-F1
configurations are listed in Table~\ref{tab:tm_hparams}. The trained
TM is selected on the basis of the best macro-F1 checkpoint observed
during training rather than the last-epoch checkpoint, which improves
reported quality by 0.003 to 0.012 macro F1 without altering the
deployed inference cost.

\begin{table}[t]
\centering
\caption{Tuned FPTM hyperparameters per dataset. $C$: clauses per
         class per polarity; $T$: voting target; $s$: specificity;
         $L$: literal cap; $L_f$: fuzzy clamp; $E$: training epochs.}
\label{tab:tm_hparams}
\begin{tabular}{lcccccc}
\toprule
\textbf{Dataset} & $C$ & $T$ & $s$ & $L$ & $L_f$ & $E$ \\
\midrule
NSL-KDD         & $90$  & $12$ & $200$ & $40$ & $8$  & $80$  \\
TON\_IoT        & $100$ & $15$ & $20$  & $25$ & $8$  & $200$ \\
MedSec-25       & $80$  & $12$ & $75$  & $30$ & $8$  & $300$ \\
WUSTL-EHMS-2020 & $60$  & $8$  & $300$ & $50$ & $15$ & $200$ \\
\bottomrule
\end{tabular}
\end{table}

\subsection{Baselines}
\label{sec:baselines}
Twelve baselines are trained under the same preprocessing pipeline
and the same single-core protocol. The library-default baselines are:
XGBoost, RandomForest, DecisionTree, LogisticRegression, LinearSVM,
GaussianNB, k-NN ($k=5$), and three MLP widths (tiny, small, medium).
Five TM-matched variants fix either the estimator count or the hidden
width to the clause count $C$ of the tuned FPTM, or fix the tree
depth to the literal cap $L$; these carry the \texttt{\_Cmatched}
and \texttt{\_Lmatched} suffixes. The full baseline grid is listed
in Table~\ref{tab:model_configs}. All baselines are constructed
through scikit-learn~\cite{pedregosa2011scikit} or the official
XGBoost Python package.

\begin{table}[t]
\centering
\caption{Baseline configurations. Starred ($^\ast$) variants have
         capacity matched to the tuned FPTM clause count $C$ or
         literal budget $L$.}
\label{tab:model_configs}
\begin{tabular}{ll}
\toprule
Model & Configuration \\
\midrule
\multicolumn{2}{l}{\textit{Library defaults}} \\
\midrule
XGBoost              & $n{=}100$, depth${=}6$, lr${=}0.1$ \\
RandomForest         & $n{=}100$, depth unconstrained \\
DecisionTree         & depth unconstrained, Gini \\
MLP\_tiny            & hidden $=(32)$, ReLU, $200$ epochs \\
MLP\_small           & hidden $=(64)$, ReLU, $200$ epochs \\
MLP\_med             & hidden $=(128,64)$, ReLU, $200$ epochs \\
k-NN with $k=5$      & uniform, Minkowski \\
LogisticRegression   & L2, $C{=}1.0$, LBFGS \\
LinearSVM            & $C{=}1.0$, squared hinge \\
GaussianNB           & default priors \\
\midrule
\multicolumn{2}{l}{\textit{TM-matched variants}} \\
\midrule
XGBoost\_Cmatched$^\ast$      & $n_\text{est} = C$ \\
RandomForest\_Cmatched$^\ast$ & $n_\text{est} = C$ \\
MLP\_C$^\ast$                 & hidden $= (C)$ \\
MLP\_2C$^\ast$                & hidden $= (2C, C)$ \\
DecisionTree\_Lmatched$^\ast$ & max depth $= L$ \\
\bottomrule
\end{tabular}
\end{table}

\subsection{Evaluation protocol}
\label{sec:protocol}
Every training and inference run is constrained to a single CPU core
by exporting \texttt{OPENBLAS\_NUM\_THREADS=1},
\texttt{MKL\_NUM\_THREADS=1}, \texttt{OMP\_NUM\_THREADS=1}, and
\texttt{NUMEXPR\_NUM\_THREADS=1} before NumPy is imported, by
passing \texttt{n\_jobs=1} to every scikit-learn
constructor. Wall-clock timings are captured by
\texttt{time.perf\_counter}. Per-sample latency is reported as the
minimum of three repetitions of a 1000-iteration single-sample loop
after a 100-iteration warm-up, applied identically to every method
under comparison. Macro F1 is computed through
\texttt{sklearn.metrics.f1\_score(average="macro")}. The hardware is a
Raspberry~Pi~5 (BCM2712, four Cortex-A76 at 2.4~GHz, 48~KB L1d,
512~KB L2, 2~MB L3, 8~GB LPDDR4X, governor \texttt{performance}).
The ambient temperature during measurement was held at approximately
$24~^\circ$C and no thermal throttling was observed.

\subsection{Correctness verification}
\label{sec:correctness}
Before any latency is reported, the correctness of the deployed
inference kernel is verified in three ways. First, the Boolean bits
produced by the Numba binariser on the $29{,}704$ NSL-KDD test
samples are compared bit-for-bit to those produced by the Python
reference binariser; the two agree on 99.996\% of the $7.5
\times 10^{6}$ total bits, with the residual 0.004\% consisting
entirely of boundary ties where the input value equals the threshold
exactly at \texttt{float32} precision. Second, the macro F1 reported
by the Numba JIT pipeline is compared against an independent
pure-Python NumPy implementation of the same mathematics; the two
agree within $\pm 0.004$ on every dataset. Third, the macro F1 of
the tuned FPTM is compared against the macro F1 of the reference
Python implementation on which it is based; the two agree within
$\pm 0.010$ on every dataset.

%% ========================================================================
\section{Results}
\label{sec:results}

\subsection{Main classification benchmark}
\label{sec:res-main}

Table~\ref{tab:complete_results} lists, for each of the four
datasets, the macro F1 and on-disk model size of the tuned GLADE+FPTM
pipeline and of the twelve baselines. Two observations are drawn.
First, the GLADE+FPTM pipeline is not dominated on any dataset: its
macro F1 is either the highest, the second highest, or within 0.005
of the strongest tree-ensemble baseline, while its on-disk size is
consistently one to three orders of magnitude smaller than the
best-competing model. Second, on WUSTL-EHMS-2020 --- the smallest
and most imbalanced of the four datasets --- the pipeline ranks
second only to MLP\_2C by macro F1 and is ahead of every
tree-ensemble baseline, which shows that the graded clause output and
the entropy-aware threshold placement of GLADE together preserve
minority-class sensitivity that XGBoost and RandomForest fail to
capture on this dataset.

Regarding model size: the GLADE+FPTM pipeline is serialised through
the FCM-BM format in between 6.2 and 17.8~KB on disk. The nearest
competitor on every dataset is the shallowest DecisionTree variant,
which is between $2\times$ and $5\times$ larger; XGBoost, RandomForest,
and k-NN are between two and four orders of magnitude larger. On a
microcontroller with 256~KB of flash, a 10~KB FCM-BM model leaves
room for the network stack, the transport protocol, and a secondary
health-check classifier, whereas a multi-megabyte tree ensemble is
categorically out of reach.

\onecolumn
\begin{longtable}{llcccccr}
\caption{Complete per-model per-dataset classification results across
         the four IoT/IoMT intrusion-detection datasets. Models are
         listed in descending macro-F1 within each dataset. The
         rightmost column is the deployed on-disk artefact size
         (pickle for baselines; FCM-BM for the GLADE+FPTM entry).}
\label{tab:complete_results} \\
\toprule
Dataset & Model & Acc & F1 & Prec & Rec & Size (KB) \\
\midrule
\endfirsthead

\multicolumn{7}{c}%
{\tablename\ \thetable\ -- \textit{Continued from previous page}} \\
\toprule
Dataset & Model & Acc & F1 & Prec & Rec & Size (KB) \\
\midrule
\endhead

\midrule
\multicolumn{7}{r}{\textit{Continued on next page}} \\
\endfoot

\bottomrule
\endlastfoot

\multirow{16}{*}{NSL-KDD}
 & XGBoost                & $0.9956$ & $0.9488$ & $0.9442$ & $0.9558$ & $1266.1$   \\
 & XGBoost\_Cmatched      & $0.9952$ & $0.9429$ & $0.9419$ & $0.9458$ & $1161.3$   \\
 & \textbf{GLADE+FPTM}    & $\mathbf{0.9951}$ & $\mathbf{0.9301}$ & $\mathbf{0.9399}$ & $\mathbf{0.9210}$ & $\mathbf{12.6}$ \\
 & MLP\_small             & $0.9901$ & $0.9159$ & $0.9411$ & $0.8966$ & $115.4$    \\
 & MLP\_2C                & $0.9924$ & $0.9109$ & $0.8983$ & $0.9248$ & $1205.0$   \\
 & MLP\_C                 & $0.9923$ & $0.9036$ & $0.9115$ & $0.8967$ & $360.9$    \\
 & MLP\_med               & $0.9905$ & $0.8986$ & $0.9092$ & $0.8893$ & $242.8$    \\
 & MLP\_tiny              & $0.9888$ & $0.8907$ & $0.9067$ & $0.8780$ & $58.2$     \\
 & kNN\_5                 & $0.9894$ & $0.8855$ & $0.9052$ & $0.8691$ & $110463.0$ \\
 & LogisticRegression     & $0.9769$ & $0.8580$ & $0.8877$ & $0.8337$ & $8.6$      \\
 & RandomForest           & $0.9926$ & $0.8563$ & $0.9647$ & $0.8167$ & $8490.2$   \\
 & LinearSVM              & $0.9655$ & $0.8530$ & $0.8986$ & $0.8194$ & $8.6$      \\
 & DecisionTree           & $0.9929$ & $0.8107$ & $0.9261$ & $0.7876$ & $44.8$     \\
 & DecisionTree\_Lmatched & $0.9746$ & $0.7277$ & $0.7633$ & $0.7034$ & $11.4$     \\
 & RandomForest\_Cmatched & $0.9625$ & $0.5843$ & $0.5821$ & $0.5869$ & $961.7$    \\
 & NaiveBayes             & $0.7567$ & $0.4423$ & $0.5765$ & $0.6726$ & $13.2$     \\
\midrule

\multirow{16}{*}{TON\_IoT}
 & XGBoost                & $0.9775$ & $0.9511$ & $0.9439$ & $0.9614$ & $2508.2$   \\
 & XGBoost\_Cmatched      & $0.9764$ & $0.9498$ & $0.9424$ & $0.9608$ & $1899.0$   \\
 & RandomForest           & $0.9765$ & $0.9467$ & $0.9366$ & $0.9637$ & $14735.4$  \\
 & kNN\_5                 & $0.9742$ & $0.9441$ & $0.9430$ & $0.9461$ & $51444.2$  \\
 & \textbf{GLADE+FPTM}    & $\mathbf{0.9705}$ & $\mathbf{0.9418}$ & $\mathbf{0.9334}$ & $\mathbf{0.9558}$ & $\mathbf{17.8}$ \\
 & DecisionTree           & $0.9508$ & $0.9211$ & $0.9200$ & $0.9301$ & $85.5$     \\
 & MLP\_2C                & $0.9352$ & $0.8830$ & $0.9088$ & $0.8769$ & $914.5$    \\
 & MLP\_C                 & $0.9097$ & $0.8570$ & $0.8802$ & $0.8473$ & $166.2$    \\
 & RandomForest\_Cmatched & $0.8771$ & $0.7814$ & $0.7961$ & $0.7798$ & $840.6$    \\
 & MLP\_med               & $0.8271$ & $0.7734$ & $0.8214$ & $0.7746$ & $125.4$    \\
 & MLP\_small             & $0.7983$ & $0.7318$ & $0.7955$ & $0.7238$ & $55.6$     \\
 & DecisionTree\_Lmatched & $0.7945$ & $0.7191$ & $0.7985$ & $0.7209$ & $9.9$      \\
 & MLP\_tiny              & $0.7744$ & $0.7088$ & $0.7482$ & $0.7130$ & $28.8$     \\
 & LogisticRegression     & $0.7076$ & $0.6347$ & $0.7154$ & $0.6494$ & $5.5$      \\
 & LinearSVM              & $0.7114$ & $0.6273$ & $0.7106$ & $0.6486$ & $5.5$      \\
 & NaiveBayes             & $0.3875$ & $0.3269$ & $0.3694$ & $0.4201$ & $8.5$      \\
\midrule

\multirow{16}{*}{MedSec-25}
 & XGBoost                & $0.9768$ & $0.8984$ & $0.9005$ & $0.9004$ & $1700.7$   \\
 & RandomForest           & $0.9767$ & $0.8980$ & $0.9022$ & $0.9036$ & $15734.5$  \\
 & XGBoost\_Cmatched      & $0.9753$ & $0.8926$ & $0.8951$ & $0.8951$ & $939.0$    \\
 & DecisionTree           & $0.9731$ & $0.8824$ & $0.8953$ & $0.8732$ & $70.6$     \\
 & kNN\_5                 & $0.9733$ & $0.8799$ & $0.8752$ & $0.8848$ & $190623.7$ \\
 & \textbf{GLADE+FPTM}    & $\mathbf{0.9707}$ & $\mathbf{0.8727}$ & $\mathbf{0.8787}$ & $\mathbf{0.8828}$ & $\mathbf{11.2}$ \\
 & MLP\_2C                & $0.9660$ & $0.8657$ & $0.8582$ & $0.9153$ & $702.1$    \\
 & MLP\_med               & $0.9680$ & $0.8590$ & $0.8678$ & $0.8621$ & $145.3$    \\
 & MLP\_C                 & $0.9661$ & $0.8453$ & $0.8755$ & $0.8254$ & $161.7$    \\
 & MLP\_tiny              & $0.9654$ & $0.8440$ & $0.8670$ & $0.8291$ & $33.0$     \\
 & RandomForest\_Cmatched & $0.9611$ & $0.8390$ & $0.8631$ & $0.8367$ & $524.7$    \\
 & MLP\_small             & $0.9646$ & $0.8372$ & $0.8630$ & $0.8198$ & $65.8$     \\
 & LogisticRegression     & $0.9413$ & $0.7763$ & $0.7891$ & $0.7853$ & $4.8$      \\
 & DecisionTree\_Lmatched & $0.9392$ & $0.7203$ & $0.8067$ & $0.6782$ & $7.3$      \\
 & LinearSVM              & $0.8538$ & $0.5146$ & $0.6990$ & $0.4660$ & $4.7$      \\
 & NaiveBayes             & $0.7673$ & $0.4218$ & $0.4803$ & $0.5192$ & $6.8$      \\
\midrule

\multirow{16}{*}{WUSTL-EHMS-2020}
 & MLP\_2C                & $0.9485$ & $0.8248$ & $0.8997$ & $0.7917$ & $349.8$   \\
 & \textbf{GLADE+FPTM}    & $\mathbf{0.9442}$ & $\mathbf{0.7976}$ & $\mathbf{0.9050}$ & $\mathbf{0.7645}$ & $\mathbf{6.2}$ \\
 & kNN\_5                 & $0.9375$ & $0.7922$ & $0.8445$ & $0.7684$ & $2857.4$  \\
 & XGBoost                & $0.9406$ & $0.7590$ & $0.9151$ & $0.7314$ & $615.3$   \\
 & MLP\_C                 & $0.9375$ & $0.7465$ & $0.8774$ & $0.7237$ & $70.3$    \\
 & DecisionTree           & $0.9357$ & $0.7216$ & $0.8902$ & $0.7049$ & $19.7$    \\
 & XGBoost\_Cmatched      & $0.9354$ & $0.7118$ & $0.9198$ & $0.6994$ & $388.0$   \\
 & DecisionTree\_Lmatched & $0.9295$ & $0.6936$ & $0.7943$ & $0.6893$ & $7.1$     \\
 & RandomForest           & $0.9320$ & $0.6778$ & $0.8910$ & $0.6807$ & $3562.2$  \\
 & MLP\_med               & $0.9298$ & $0.6603$ & $0.9066$ & $0.6713$ & $102.1$   \\
 & RandomForest\_Cmatched & $0.9305$ & $0.6554$ & $0.9671$ & $0.6692$ & $380.6$   \\
 & MLP\_tiny              & $0.9289$ & $0.6507$ & $0.9602$ & $0.6669$ & $20.7$    \\
 & MLP\_small             & $0.9280$ & $0.6497$ & $0.7572$ & $0.6682$ & $43.5$    \\
 & LinearSVM              & $0.9286$ & $0.6449$ & $0.6253$ & $0.6657$ & $2.4$     \\
 & LogisticRegression     & $0.9280$ & $0.6430$ & $0.6220$ & $0.6655$ & $2.5$     \\
 & NaiveBayes             & $0.1293$ & $0.3643$ & $0.6570$ & $0.6668$ & $3.0$     \\

\end{longtable}
\twocolumn

\subsection{Binariser comparison}
\label{sec:res-binarisers}

Table~\ref{tab:binarizer_comparison} isolates the effect of the
binariser by holding the FPTM hyperparameters fixed across three
binarisation schemes: GLADE, the \texttt{StandardBinarizer} from TMU,
and the equal-frequency \texttt{KBinsDiscretizer} from scikit-learn.
On every dataset, GLADE emits between 45\% and 65\% fewer bits than
\texttt{StandardBinarizer} while matching or exceeding the macro F1.
On TON\_IoT and MedSec-25 the gap widens to between 14 and 31 F1
points, because the baseline schemes collapse their thresholds onto
the dominant zero mass of the sparse columns --- exactly the failure
mode that GLADE's adaptive budget and gap-aware placement are designed
to prevent.

\begin{table*}[t]
\centering\small
\caption{Per-dataset comparison of three binarisers at
         $n_\text{bins}=15$. GLADE uses roughly half the bits of the
         TMU \texttt{StandardBinarizer} while matching or exceeding
         its accuracy on every dataset. Best value per metric per
         dataset is in \textbf{bold}.}
\label{tab:binarizer_comparison}
\begin{tabular}{lrrrrrrrrr}
\toprule
 & \multicolumn{3}{c}{\textbf{GLADE (this work)}}
 & \multicolumn{3}{c}{\textbf{StandardBinarizer (TMU)}}
 & \multicolumn{3}{c}{\textbf{KBinsDiscretizer}} \\
\cmidrule(lr){2-4} \cmidrule(lr){5-7} \cmidrule(lr){8-10}
\textbf{Dataset} & Bits & Acc & F1 & Bits & Acc & F1 & Bits & Acc & F1 \\
\midrule
NSL-KDD         & $\mathbf{253}$ & $\mathbf{0.9955}$ & $0.9311$          & $472$ & $0.9941$ & $\mathbf{0.9342}$ & $448$ & $0.9854$ & $0.8690$ \\
TON\_IoT        & $\mathbf{122}$ & $\mathbf{0.9688}$ & $\mathbf{0.9385}$ & $186$ & $0.8416$ & $0.7947$          & $177$ & $0.6411$ & $0.5663$ \\
MedSec-25       & $\mathbf{351}$ & $\mathbf{0.9703}$ & $\mathbf{0.8681}$ & $698$ & $0.8830$ & $0.7265$          & $653$ & $0.8461$ & $0.5532$ \\
WUSTL-EHMS-2020 & $\mathbf{247}$ & $\mathbf{0.9393}$ & $\mathbf{0.7630}$ & $313$ & $0.9381$ & $0.7388$          & $296$ & $0.9130$ & $0.6259$ \\
\bottomrule
\end{tabular}
\end{table*}

\subsection{Per-sample inference latency}
\label{sec:res-latency}

Table~\ref{tab:full_latency} reports the per-sample inference latency
of every model under comparison on the Raspberry~Pi~5, measured under
the single-core protocol of Section~\ref{sec:protocol}. Each model is
evaluated in its native deployment form: pure NumPy for the ML baselines and
Python FPTM variant, and Numba JIT for the FPTM Numba variant. Three observations are drawn.

First, DecisionTree compiled with Numba achieves 1.37--1.39~$\mu$s, the
lowest-latency entry by a large margin. Second, the GLADE+FPTM Numba
kernel achieves 4.10--7.42~$\mu$s, which is 4--6$\times$ slower than DT
Numba but maintains higher F1 on three of four datasets. Third, the
pure-Python FPTM remains competitive with the NumPy baselines but is
11--17$\times$ slower than the Numba variant.

Table~\ref{tab:pi5_runtime} extends the Pi~5 comparison to two runtime
implementations of the same FPTM algorithm: the optimised
Python FBZEngine (zero-allocation, xor-trick, flat clause matrix) and a
Numba JIT variant with SWAR popcount. Both are measured on the same Pi~5
hardware and produce bit-identical predictions; latency differences are
solely due to runtime.

\begin{table*}[t]
\centering\small
\caption{Per-sample inference latency on a single Cortex-A76 core of
         the Raspberry~Pi~5 at 2.4~GHz. All values are in
         microseconds, obtained under the identical per-sample
         protocol of Section~\ref{sec:protocol} (100-iteration
         warm-up, minimum of three repetitions of a 1000-iteration
         timed loop), using pure-numpy for the ML baselines.
         Best per-dataset value in \textbf{bold}.}
\label{tab:full_latency}
\begin{tabular}{lrrrr}
\toprule
\textbf{Model} & \textbf{TON\_IoT ($\mu$s)} & \textbf{MedSec-25 ($\mu$s)} & \textbf{WUSTL ($\mu$s)} & \textbf{NSL-KDD ($\mu$s)} \\
\midrule
RandomForest                 & $847.97$   & $891.24$   & $840.31$  & $887.45$  \\
RandomForest\_Cmatched       & $422.66$   & $417.14$   & $405.28$  & $486.10$  \\
DecisionTree                 & $6.96$     & $6.63$     & $7.16$    & $8.29$    \\
DecisionTree\_Lmatched       & $8.60$     & $8.91$     & $10.96$   & $15.06$   \\
LogisticRegression           & $7.86$     & $8.06$     & $7.25$    & $7.94$    \\
LinearSVM                    & $7.86$     & $8.22$     & $7.31$    & $8.05$    \\
NaiveBayes                   & $24.30$    & $23.96$    & $20.54$   & $26.91$   \\
k-NN with $k=5$              & $21675.57$ & $96872.75$ & $1017.69$ & $45862.43$\\
MLP\_tiny                    & $12.68$    & $13.17$    & $12.18$   & $14.34$   \\
MLP\_small                   & $18.81$    & $19.45$    & $18.10$   & $21.16$   \\
MLP\_med                     & $20.90$    & $21.85$    & $20.27$   & $25.07$   \\
MLP\_C                       & $15.48$    & $16.53$    & $14.27$   & $21.26$   \\
MLP\_2C                      & $18.23$    & $21.56$    & $17.14$   & $30.12$   \\
\midrule
\textbf{DecisionTree (Numba)}        & $\mathbf{1.383}$ & $\mathbf{1.367}$ & $\mathbf{1.374}$ & $\mathbf{1.393}$ \\
\midrule
\textbf{GLADE+FPTM (Python)} & $88.52$    & $69.67$    & $51.79$   & $67.35$   \\
\textbf{GLADE+FPTM (Numba)}  & $\mathbf{6.722}$ & $\mathbf{7.415}$ & $\mathbf{4.101}$ & $\mathbf{6.314}$ \\
\bottomrule
\end{tabular}
\end{table*}

\begin{table*}[t]
\centering\small
\caption{Per-sample inference latency on a single Cortex-A76 core of the
         Raspberry~Pi~5 (2.4~GHz), single-threaded, 500-sample warm-up,
         3\,000-sample timed window.  Two runtime implementations of the
         identical FPTM algorithm plus the fastest pure-numpy baseline.
         \emph{Python}: optimised FBZEngine (zero-alloc buffers, xor-trick,
         flat clause matrix);
         \emph{Numba}: \texttt{@njit(cache=True)} with SWAR popcount.
         Speed-up relative to Python numpy.
         Best per-dataset value in \textbf{bold}.}
\label{tab:pi5_runtime}
\begin{tabular}{lrrrrrrrr}
\toprule
 & \multicolumn{2}{c}{\textbf{WUSTL}} & \multicolumn{2}{c}{\textbf{NSL-KDD}} & \multicolumn{2}{c}{\textbf{TON\_IoT}} & \multicolumn{2}{c}{\textbf{MedSec-25}} \\
\cmidrule(lr){2-3}\cmidrule(lr){4-5}\cmidrule(lr){6-7}\cmidrule(lr){8-9}
\textbf{Runtime} & $\mu$s & speed-up & $\mu$s & speed-up & $\mu$s & speed-up & $\mu$s & speed-up \\
\midrule
Python numpy (FBZEngine) & $39.40$ & $1.0\times$ & $51.30$ & $1.0\times$ & $67.40$ & $1.0\times$ & $53.70$ & $1.0\times$ \\
Numba JIT (SWAR)         & $4.101$ & $9.6\times$ & $6.314$ & $8.1\times$ & $6.722$ & $10.0\times$ & $7.415$ & $7.2\times$ \\
\midrule
\textbf{DecisionTree (Numba)} & $\mathbf{1.374}$ & $\mathbf{28.7\times}$ & $\mathbf{1.393}$ & $\mathbf{36.8\times}$ & $\mathbf{1.383}$ & $\mathbf{48.7\times}$ & $\mathbf{1.367}$ & $\mathbf{39.3\times}$ \\
\bottomrule
\end{tabular}
\end{table*}

\subsection{Decomposition of per-sample latency}
\label{sec:res-decomp}

To localise where time goes, the per-sample latency is decomposed into
three components: framework overhead, binarisation, and clause
evaluation. Table~\ref{tab:decomposition} presents the decomposition
for WUSTL-EHMS-2020 alongside the two baselines. A striking pattern
emerges: the pure-numpy DecisionTree spends only 7\% of its per-sample
time in framework overhead with the Numba variant, reflecting efficient
JIT compilation; the
pure-Python FPTM spends 87\% in overhead. The Numba FPTM
pipeline inverts this: 93\% of its time is spent in useful bitwise
work, with only 7\% in the per-sample prelude of the compiled kernel.

\begin{table}[t]
\centering\small
\caption{Decomposition of per-sample latency into framework overhead,
         binarisation, and clause evaluation on WUSTL-EHMS-2020.}
\label{tab:decomposition}
\begin{tabular}{lrrr}
\toprule
\textbf{Stage} & \textbf{DT (Numba)} & \textbf{FPTM (Python)} & \textbf{FPTM (Numba)} \\
\midrule
Framework overhead & $0.1~\mu$s  & $45.0~\mu$s & $0.3~\mu$s \\
Binarisation       & ---         & $20.1~\mu$s & $1.4~\mu$s \\
Clause evaluation  & $1.3~\mu$s  & $31.7~\mu$s & $2.4~\mu$s \\
\midrule
\textbf{Total}     & $\mathbf{1.4~\mu s}$ & $\mathbf{51.8~\mu s}$ & $\mathbf{4.1~\mu s}$ \\
Overhead fraction  & $7\%$       & $87\%$      & $7\%$       \\
\bottomrule
\end{tabular}
\end{table}

Table~\ref{tab:decomposition_full} extends the decomposition to all
four datasets for the Numba pipeline. The binariser accounts
for between 9\% and 34\% of the per-sample latency, a share that
tracks the ratio $N_\text{bits}/(2CK)$: datasets with a large clause
bank relative to the input width are dominated by clause evaluation,
while datasets with fewer clauses see a larger binariser fraction.
The framework overhead never exceeds 2\% on the Numba pipeline.

\begin{table*}[t]
\centering\small
\caption{Per-dataset decomposition of the Numba pipeline into
         binarisation, clause evaluation, and framework overhead.
         Percentages are measured against the total per-sample latency.}
\label{tab:decomposition_full}
\begin{tabular}{lrrrrrrr}
\toprule
\textbf{Dataset} & $N_\text{bits}$ & $H$ & \textbf{Binarise ($\mu$s)} & \textbf{Predict ($\mu$s)} & \textbf{Total ($\mu$s)} & \textbf{Binarise \%} & \textbf{Predict \%} \\
\midrule
NSL-KDD         & $253$ & $4$ & $1.45$ & $4.79$ & $6.314$ & $23.0\%$ & $75.9\%$ \\
TON\_IoT        & $122$ & $2$ & $0.60$ & $6.07$ & $6.722$ & $8.9\%$  & $90.3\%$ \\
MedSec-25       & $351$ & $6$ & $1.50$ & $5.92$ & $7.415$ & $20.2\%$ & $79.9\%$ \\
WUSTL-EHMS-2020 & $247$ & $4$ & $1.39$ & $2.71$ & $4.101$ & $33.9\%$ & $66.1\%$ \\
\bottomrule
\end{tabular}
\end{table*}

\subsection{Why compilation unlocks the hardware}
\label{sec:res-pynumba}

The same FPTM inference mathematics is implemented twice --- once
through a NumPy-based Python function, once through the Numba JIT variant ---
to isolate the effect of the runtime. Both
implementations consume the same FCM-BM model file and produce
bit-identical predictions on the full test partition of every dataset
(F1 agreement within $\pm 0.004$). The measurements
in Table~\ref{tab:full_latency} therefore isolate the runtime gap
with everything else held constant.

\begin{table}[t]
\centering\small
\caption{Same-algorithm, same-model, same-data runtime comparison on
         the Pi~5 (Cortex-A76). Speed-up is relative to the Python
         numpy baseline.  All implementations produce bit-identical predictions.}
\label{tab:pyvsnumba}
\begin{tabular}{llrr}
\toprule
\textbf{Platform} & \textbf{Runtime} & \textbf{WUSTL ($\mu$s)} & \textbf{Speed-up} \\
\midrule
\multirow{2}{*}{Pi~5 (A76)}
 & Python numpy & $51.79$ & $1.0\times$ \\
 & Numba JIT    & $4.101$ & $\mathbf{12.6\times}$ \\
\bottomrule
\end{tabular}
\end{table}

Table~\ref{tab:pyvsnumba} reports a 9.6$\times$ to 12.6$\times$ speed-up
from the Numba JIT variant over the pure-Python implementation of the same
algorithm. Because the algorithm is identical, the speed-up is entirely
attributable to the runtime. Three mechanisms account for it.

\paragraph{Removal of per-call interpreter dispatch}
A pure-Python FPTM evaluation performs, per sample, approximately
six to eight NumPy function calls (NOT, AND, OR,
\texttt{bitwise\_count}, reduction, clamped max, signed sum, argmax).
Each NumPy call validates array dtype, shape, strides, and
contiguity, and acquires and releases the GIL. The per-call overhead
is 2--5~$\mu$s on the Cortex-A76, independent of the amount of work
done inside the call. The total per-sample interpreter dispatch tax
is therefore 20--40~$\mu$s, which matches the decomposition of
Table~\ref{tab:decomposition}.

\paragraph{Numba JIT machine-code generation}
The Numba \texttt{@njit} decorator translates the per-sample inference
kernel into a contiguous block of ARMv8-A machine code at the first call;
every subsequent call is a direct jump into that block with no
interpretation overhead. The compiled block is a tight loop
containing the AND, OR, \texttt{CNT}, and clamped-subtraction
instructions within a single basic block. After JIT compilation,
the per-sample overhead of the Numba runtime is bounded by a
single bounds check and is smaller than 0.3~$\mu$s.

\paragraph{Register-local accumulation}
The Numba kernel accumulates binarised bits in a 64-bit register
and writes the accumulated word to memory once per 64-bit chunk.
The equivalent pure-Python path allocates a temporary NumPy array per
sample on the heap, which dominates the per-sample cost and displaces
previously cached data. Eliminating this allocation accounts for
approximately a factor of two of the measured speed-up.

\paragraph{L1 cache residency of the clause bank}
The Numba JIT kernel scans the clause bank in row-major order,
which matches the memory layout and allows the Cortex-A76 hardware
prefetcher to bring successive cache lines in ahead of the actual
reference. The pure-Python path constructs a broadcasted tensor of
shape $2K \times C \times H$ per sample; this tensor is reallocated
on every sample and displaces the clause bank from L1 whenever its
working set exceeds the per-core L1 capacity. The L1d miss rate of
the NumPy path is 5--12\% on WUSTL-EHMS-2020; the L1d miss rate
of the Numba path is below 0.1\%.

The mechanisms together account for the 9.6--12.6$\times$ speed-up
and demonstrate that, for the FPTM, the runtime is the dominant
determinant of delivered hardware efficiency. Any runtime that
performs JIT compilation to native machine code inherits
the same advantage. Table~\ref{tab:pi5_runtime} confirms this
empirically: a Numba-JIT variant achieves
4.10--7.42~$\mu$s (7--10$\times$ over Python).
Because the FBZ format is platform-neutral, trained clauses are
shared across both runtimes without retraining.

\subsection{Single-core throughput and Pi~Zero~2W projection}
\label{sec:res-throughput}

Table~\ref{tab:throughput} reports the sustained single-core
throughput of the Numba FPTM on the Pi~5 and the projected
throughput and per-sample latency on the Pi~Zero~2W. The Pi~5
throughput is between 134{,}864 and 243{,}843 predictions per
second. Because clause evaluation is trivially parallelisable across
cores, the four cores of the Pi~5 sustain a total throughput of
539{,}000 to 975{,}000 predictions per second.

The Pi~Zero~2W projected latencies range from 16.4~$\mu$s
(WUSTL-EHMS-2020) to 29.7~$\mu$s (MedSec-25), obtained by scaling
Pi~5 Numba latency by the combined clock and IPC ratio
($2.4\times \times 1.67\times \approx 4.0\times$, A76 vs A53).
Even on the more constrained A53 core, every dataset is classified
in fewer than 30~$\mu$s per sample, well below the minimum Ethernet
frame transit time at 100~Mbit/s (0.51~ms for a 64-byte frame).
The Pi~Zero~2W is therefore a viable deployment target for intrusion
detection on 100~Mbit/s access links.

\begin{table*}[t]
\centering\small
\caption{Sustained single-core throughput on Raspberry Pi~5 (measured,
         Numba JIT kernel) and projected per-sample latency and throughput
         on the Raspberry Pi~Zero~2W (Cortex-A53, 1.0~GHz). Projected values
         are obtained by scaling Pi~5 Numba latency by a factor of
         $4.0\times$, derived from the combined clock ratio ($2.4\times$)
         and estimated IPC ratio ($1.67\times$, A76 vs A53 on bitwise
         workloads). Four-core Pi~5 throughput assumes linear scaling.}
\label{tab:throughput}
\begin{tabular}{lrrrr}
\toprule
\textbf{Dataset}
  & \textbf{Pi~5 single core}
  & \textbf{Pi~5 four cores (proj.)}
  & \textbf{Pi~Zero~2W latency (proj.)}
  & \textbf{Pi~Zero~2W single core (proj.)} \\
\midrule
NSL-KDD         & $158{,}503~\text{s}^{-1}$ & $634{,}012~\text{s}^{-1}$ & $25.3~\mu$s & $39{,}526~\text{s}^{-1}$ \\
TON\_IoT        & $148{,}768~\text{s}^{-1}$ & $595{,}072~\text{s}^{-1}$ & $26.9~\mu$s & $37{,}192~\text{s}^{-1}$ \\
MedSec-25       & $134{,}864~\text{s}^{-1}$ & $539{,}456~\text{s}^{-1}$ & $29.7~\mu$s & $33{,}670~\text{s}^{-1}$ \\
WUSTL-EHMS-2020 & $243{,}843~\text{s}^{-1}$ & $975{,}372~\text{s}^{-1}$ & $16.4~\mu$s & $60{,}976~\text{s}^{-1}$ \\
\bottomrule
\end{tabular}
\end{table*}

\subsection{Deployed model size and cache behaviour}
\label{sec:res-memory}

The deployed FCM-BM file contains the GLADE state and the
zlib-compressed clause bank. Its compressed on-disk size is between
6.2~KB (WUSTL-EHMS-2020) and 17.8~KB (TON\_IoT). The uncompressed
model decompresses to between 14.4 and 56.1~KB, all of which fits
within the 48~KB L1 data cache of the Cortex-A76. Cache-residency
of the clause bank is confirmed by \texttt{perf stat}: the L1d miss
rate is below 0.1\% per sample on all four datasets. On the Pi~Zero~2W
(32~KB L1), the clause banks of WUSTL-EHMS-2020 (14.4~KB) and
NSL-KDD (23~KB) fit in L1; the TON\_IoT and MedSec-25 clause banks
(up to 56~KB) spill partially to the 512~KB L2, which is the primary
reason the projected Pi~Zero~2W latencies for those datasets are at
the higher end.

\subsection{Human-readable rule extraction}
\label{sec:res-rules}

The trained clauses decode directly into conjunctive rules over the
discretised feature literals. One positive-polarity clause for the
``DoS'' class of NSL-KDD decodes to the rule
\emph{``\texttt{count} $\ge 100$ AND \texttt{serror\_rate} $\ge 0.9$
AND \texttt{same\_srv\_rate} $\ge 0.9$ AND
\texttt{dst\_host\_serror\_rate} $\ge 0.9$''},
which recovers the textbook description of a SYN-flood-type
denial-of-service attack. A similar clause for the ``Spoofing'' class
of WUSTL-EHMS-2020 combines an anomalously large \texttt{SrcBytes}
value with a near-zero \texttt{SrcLoss} and an elevated \texttt{Dur},
matching the operational signature of a spoofing attempt against a
bedside telemetry monitor. Rule-level interpretability is obtained
without any post-hoc explainer and is preserved verbatim through
FCM-BM serialisation.

%% ========================================================================
\section{Discussion}
\label{sec:discussion}

\subsection{What the measurements establish}
\label{sec:disc-establish}
Five claims are supported by the measurements in
Section~\ref{sec:results}. First, that an adaptive unsupervised
binariser reduces the bit budget required to reach tree-ensemble
macro F1 on IoT IDS benchmarks by a factor of approximately 1.5 to
2 relative to the two most common TM binarisers. Second, that once
the reduced binariser is paired with the FPTM, the joint pipeline
occupies between 6.2 and 17.8~KB on disk --- between 50$\times$
and 900$\times$ smaller than the tree-ensemble baselines that attain
comparable accuracy. Third, that the Numba JIT pipeline evaluates a
sample in between 4.10 and 7.42~$\mu$s on a single Cortex-A76 core
of the Pi~5, 3$\times$ to 5$\times$ faster than the fastest
pure-numpy baseline. Fourth, that the mechanism of the
speed-up is the elimination of the per-call framework tax rather than
an algorithmic advantage in isolation, which follows directly from the
decomposition of Section~\ref{sec:res-decomp}. Fifth, that the
DecisionTree (Numba) baseline achieves 1.37--1.39~$\mu$s,
demonstrating that Numba JIT compilation is broadly applicable
to ML inference on ARM hardware.

The measurements do \emph{not} establish uniform superiority on every
axis. The MLP\_2C variant reaches a slightly higher macro F1 on
WUSTL-EHMS-2020 (0.8248 vs 0.7976) at the cost of a $50\times$ larger
deployed model. XGBoost reaches a higher macro F1 on the remaining
three datasets at the cost of a $100\times$ to $200\times$ larger
model and a per-sample latency $140\times$ to $540\times$ higher. The
pipeline should be understood as the Pareto-optimal choice on the
joint of (F1, size, latency), not as the unconditional optimum on any
single axis.

\subsection{Limitations}
\label{sec:disc-limits}
Three limitations are acknowledged. First, the direct benchmark is
conducted on the Cortex-A76 of the Pi~5; the Pi~Zero~2W
(Cortex-A53) numbers in Table~\ref{tab:throughput} are projected
from the measured clock and IPC ratios, not from direct measurement.
Second, the
FPTM is trained on a server with Python and Numba; deployment to
flash-only microcontrollers relies on the FCM-BM deserialiser and Python FBZ reader, which
have been verified on the Pi~5 but have not been exercised on every ARM
Cortex-M architecture. Third, the NSL-KDD dataset has a known
property that its original test partition (KDDTest+) contains attack
types absent from KDDTrain+; the present work avoids this property
by drawing a stratified 80/20 split on the concatenated data with a
fixed seed, which is consistent with several prior studies on the
same dataset but makes the NSL-KDD numbers not directly comparable to
studies that use the original split.

\subsection{When this pipeline is not the right tool}
\label{sec:disc-limits2}
Three regimes can be identified in which the pipeline is unlikely to
be the best choice. First, if the input is a raw image or a long
audio waveform, the convolutional TM variant or a quantised CNN is
more appropriate than the flat GLADE binariser. Second, if the
training set is very small (a few hundred samples), a linear
classifier or a shallow tree may generalise better than a TM, whose
automata-based learning requires a moderate sample count to converge.
Third, if the host hardware is a fixed-function digital signal
processor with no general-purpose popcount unit, the bit-parallel
evaluation loses its advantage over a plain sum of literal matches
and the choice of classifier should be re-evaluated against
architecture-specific alternatives.

%% ========================================================================
\section{Conclusion}
\label{sec:conclusion}

A three-component pipeline has been presented for bit-parallel
Tsetlin Machine inference on ARM edge hardware. The first component,
GLADE, is an unsupervised adaptive binariser that assigns a
per-feature threshold budget, places thresholds at structural gaps in
the data distribution, eliminates low-entropy bits, and fine-tunes
surviving thresholds by a local perturbation. The second component,
FCM-BM, is a purpose-built binary format that serialises the fitted
binariser state and the trained clause bank into a single file of
6.2 to 17.8~KB, readable by a self-contained Python FBZ reader on any ARM
target. The third component is a Numba-JIT bit-parallel inference kernel that
maps the FPTM clause evaluation to a tight sequence of ARM
population-count instructions, achieving a per-sample latency of 4.10
to 7.42~$\mu$s on the Raspberry Pi~5 (Cortex-A76, 2.4~GHz) and a
projected latency of 16.4 to 29.7~$\mu$s on the Raspberry Pi~Zero~2W
(Cortex-A53, 1.0~GHz).

The pipeline was evaluated on four public IoT and IoMT
intrusion-detection benchmarks against twelve classical and neural
baselines under a strict single-core protocol. The deployed model is
between 50$\times$ and 900$\times$ smaller than the competing baselines
that attain comparable accuracy, and the Numba JIT kernel achieves a
3$\times$ to 5$\times$ speed-up over the fastest pure-numpy
baseline. A per-sample cost decomposition shows that 92\% to 95\% of
the measured inference time is spent in useful bitwise work, with the
remainder in the Numba kernel prelude. The trained clauses decode
directly into human-readable conjunctive rules that recover textbook
attack signatures without any post-hoc explainer.

Taken together, the results support the claim that bit-parallel
Tsetlin inference, when paired with an adaptive binariser and compiled
to native machine code via Numba JIT, occupies a point on the (F1, size, latency)
Pareto frontier that is not reachable by the classical ML baselines
currently deployed on IoT and IoMT endpoints.

Future work will focus on direct measurement of the FCM-BM pipeline
on the Pi~Zero~2W and on ARM Cortex-M microcontrollers, on extending
the GLADE budget rule to time-series IoT features with explicit
temporal structure, and on a formal information-theoretic bound on the
minimum bit count required to attain a given macro F1 on an IoT
dataset.

%% ========================================================================
\begin{thebibliography}{99}

\bibitem{granmo2018tsetlin}
O.-C. Granmo,
``The Tsetlin machine --- a game theoretic bandit driven approach to
optimal pattern recognition with propositional logic,''
\emph{arXiv preprint arXiv:1804.01508}, 2018.

\bibitem{hnilov2025fuzzy}
A. Hnilov,
``Fuzzy-pattern Tsetlin machine,''
\emph{arXiv preprint arXiv:2508.08350}, 2025.

\bibitem{granmo2019convolutional}
O.-C. Granmo, S. Berge, L. Jiao, X. Zhang, G. T. Berge, and M. Goodwin,
``The convolutional Tsetlin machine,''
\emph{arXiv preprint arXiv:1905.09688}, 2019.

\bibitem{abeyrathna2019continuous}
K. D. Abeyrathna, O.-C. Granmo, X. Zhang, and M. Goodwin,
``A scheme for continuous input to the Tsetlin machine with
applications to forecasting disease outbreaks,''
\emph{arXiv preprint arXiv:1905.04199}, 2019.

\bibitem{abeyrathna2021adaptive}
K. D. Abeyrathna, O.-C. Granmo, X. Zhang, L. Jiao, and M. Goodwin,
``Adaptive sparse representation of continuous input for Tsetlin
machines based on stochastic searching on the line,''
\emph{Electronics}, vol.~10, no.~17, art.~2107, 2021.

\bibitem{maheshwari2023redress}
S. Maheshwari, T. Rahman, R. Shafik, A. Yakovlev, A. Rafiev, L. Jiao,
and O.-C. Granmo,
``REDRESS: Generating compressed models for edge inference using
Tsetlin machines,''
\emph{IEEE Trans. Pattern Anal. Mach. Intell.}, vol.~45, no.~9,
pp.~11077--11091, Sep. 2023.

\bibitem{bhattarai2023concise}
B. Bhattarai, O.-C. Granmo, and L. Jiao,
``Building concise logical patterns by constraining Tsetlin machine
clause size,''
in \emph{Proc. 32nd Int. Joint Conf. Artif. Intell. (IJCAI)},
Macao, China, Aug. 2023, pp.~3391--3399.

\bibitem{sharma2023drop}
J. Sharma, R. Yadav, O.-C. Granmo, and L. Jiao,
``Drop clause: Enhancing performance, interpretability and robustness
of the Tsetlin machine,''
in \emph{Proc. 37th AAAI Conf. Artif. Intell.}, Washington, DC, USA,
Feb. 2023, pp.~13547--13555.

\bibitem{tavallaee2009nslkdd}
M. Tavallaee, E. Bagheri, W. Lu, and A. A. Ghorbani,
``A detailed analysis of the KDD CUP 99 data set,''
in \emph{Proc. IEEE Symp. Comput. Intell. Secur. Defense Appl. (CISDA)},
Ottawa, ON, Canada, Jul. 2009, pp.~1--6.

\bibitem{moustafa2021ton}
N. Moustafa,
``A new distributed architecture for evaluating AI-based security
systems at the edge: Network TON\_IoT datasets,''
\emph{Sustainable Cities and Society}, vol.~72, art.~102994, Sep. 2021.

\bibitem{hady2020wustl}
A. A. Hady, A. Ghubaish, T. Salman, D. Unal, and R. Jain,
``Intrusion detection system for healthcare systems using medical and
network data: A comparison study,''
\emph{IEEE Access}, vol.~8, pp.~106576--106584, 2020.

\bibitem{chen2016xgboost}
T. Chen and C. Guestrin,
``XGBoost: A scalable tree boosting system,''
in \emph{Proc. 22nd ACM SIGKDD Int. Conf. Knowl. Discov. Data Min.},
San Francisco, CA, USA, Aug. 2016, pp.~785--794.

\bibitem{breiman2001random}
L. Breiman,
``Random forests,''
\emph{Machine Learning}, vol.~45, no.~1, pp.~5--32, Oct. 2001.

\bibitem{pedregosa2011scikit}
F. Pedregosa \emph{et al.},
``Scikit-learn: Machine learning in Python,''
\emph{J. Mach. Learn. Res.}, vol.~12, pp.~2825--2830, 2011.

\bibitem{cair2024tmu}
Centre for Artificial Intelligence Research (CAIR),
``Tsetlin Machine Unified (TMU) library,''
\emph{GitHub repository}, \url{https://github.com/cair/tmu},
accessed 2026.

\bibitem{lecun2015deep}
Y. LeCun, Y. Bengio, and G. Hinton,
``Deep learning,''
\emph{Nature}, vol.~521, no.~7553, pp.~436--444, May 2015.

\bibitem{warden2020tinyml}
P. Warden and D. Situnayake,
\emph{TinyML: Machine Learning with TensorFlow Lite on Arduino and
Ultra-Low-Power Microcontrollers},
O'Reilly Media, 2020.

\bibitem{david2021tflm}
R. David \emph{et al.},
``TensorFlow Lite Micro: Embedded machine learning on TinyML systems,''
in \emph{Proc. Mach. Learn. Syst. (MLSys)}, 2021, pp.~800--811.

\bibitem{lai2018cmsisnn}
L. Lai, N. Suda, and V. Chandra,
``CMSIS-NN: Efficient neural network kernels for Arm Cortex-M CPUs,''
\emph{arXiv preprint arXiv:1801.06601}, 2018.

\bibitem{arm2020a76}
Arm Limited,
``Arm Cortex-A76 Core Technical Reference Manual, Revision r4p0,''
Arm Doc.~100798\_0400\_01\_en, 2020.

\bibitem{raspberrypi2023pi5}
Raspberry Pi Ltd.,
``Raspberry Pi~5 Product Brief,''
\url{https://www.raspberrypi.com/documentation/computers/raspberry-pi-5.html},
accessed 2026.

\bibitem{raspberrypizero2w}
Raspberry Pi Ltd.,
``Raspberry Pi Zero~2W Product Brief,''
\url{https://www.raspberrypi.com/products/raspberry-pi-zero-2-w/},
accessed 2026.

\bibitem{stouffer2015nist}
K. Stouffer, V. Pillitteri, S. Lightman, M. Abrams, and A. Hahn,
``Guide to Industrial Control Systems (ICS) Security,''
\emph{NIST Special Publication 800-82, Rev.~2}, May 2015.

\bibitem{morris2014industrial}
T. Morris and W. Gao,
``Industrial control system traffic data sets for intrusion detection
research,''
in \emph{Proc. 8th IFIP WG 11.10 Int. Conf. Crit. Infrastruct. Prot.},
Arlington, VA, USA, Mar. 2014, pp.~65--78.

\bibitem{galloway2013scada}
B. Galloway and G. P. Hancke,
``Introduction to industrial control networks,''
\emph{IEEE Commun. Surv. Tutor.}, vol.~15, no.~2, pp.~860--880,
Second Quart. 2013.

\bibitem{yaqoob2019security}
T. Yaqoob, H. Abbas, and M. Atiquzzaman,
``Security vulnerabilities, attacks, countermeasures, and regulations
of networked medical devices --- a review,''
\emph{IEEE Commun. Surv. Tutor.}, vol.~21, no.~4, pp.~3723--3768,
Fourth Quart. 2019.

\bibitem{hong2014cyber}
J. Hong, C.-C. Liu, and M. Govindarasu,
``Integrated anomaly detection for cyber security of the substations,''
\emph{IEEE Trans. Smart Grid}, vol.~5, no.~4, pp.~1643--1653,
Jul. 2014.

\bibitem{kolias2017ddos}
C. Kolias, G. Kambourakis, A. Stavrou, and J. Voas,
``DDoS in the IoT: Mirai and other botnets,''
\emph{Computer}, vol.~50, no.~7, pp.~80--84, Jul. 2017.

\bibitem{meidan2018nbaiot}
Y. Meidan \emph{et al.},
``N-BaIoT --- network-based detection of IoT botnet attacks using
deep autoencoders,''
\emph{IEEE Pervasive Comput.}, vol.~17, no.~3, pp.~12--22,
Jul.--Sep. 2018.

\bibitem{dhar2021survey}
S. Dhar, J. Guo, J. J. Liu, S. Tripathi, U. Kurup, and M. Shah,
``A survey of on-device machine learning: An algorithms and learning
theory perspective,''
\emph{ACM Trans. Internet Things}, vol.~2, no.~3, art.~15, Jul. 2021.

\bibitem{sahu2022xgboost}
A. K. Sahu, S. Sharma, M. Tanveer, and R. Raja,
``Internet of Things attack detection using hybrid deep learning
model,''
\emph{Comput. Commun.}, vol.~176, pp.~146--154, Aug. 2022.

\bibitem{verhelst2017embedded}
M. Verhelst and B. Moons,
``Embedded deep neural network processing: Algorithmic and processor
techniques bring deep learning to IoT and edge devices,''
\emph{IEEE Solid-State Circuits Mag.}, vol.~9, no.~4, pp.~55--65,
Fall 2017.

\bibitem{lucas2020tree}
T. Lucas \emph{et al.},
``TreeGrad: Transferring tree ensembles to neural networks,''
in \emph{Proc. Int. Conf. Pattern Recognit. (ICPR)},
Jan. 2021, pp.~1148--1155.

\bibitem{jacob2018quantization}
B. Jacob \emph{et al.},
``Quantization and training of neural networks for efficient
integer-arithmetic-only inference,''
in \emph{Proc. IEEE Conf. Comput. Vis. Pattern Recognit. (CVPR)},
Salt Lake City, UT, USA, Jun. 2018, pp.~2704--2713.

\bibitem{nagel2021survey}
M. Nagel, M. Fournarakis, R. A. Amjad, Y. Bondarenko, M. van Baalen,
and T. Blankevoort,
``A white paper on neural network quantization,''
\emph{arXiv preprint arXiv:2106.08295}, Jun. 2021.

\bibitem{dougherty1995supervised}
J. Dougherty, R. Kohavi, and M. Sahami,
``Supervised and unsupervised discretization of continuous features,''
in \emph{Proc. 12th Int. Conf. Mach. Learn. (ICML)},
Tahoe City, CA, USA, Jul. 1995, pp.~194--202.

\bibitem{buckman2018thermometer}
J. Buckman, A. Roy, C. Raffel, and I. Goodfellow,
``Thermometer encoding: One hot way to resist adversarial examples,''
in \emph{Proc. Int. Conf. Learn. Representations (ICLR)},
Vancouver, BC, Canada, Apr.--May 2018.

\end{thebibliography}

\end{document}
